
\chapter{\MakeUppercase{Methodology and Testing}}

\section{Overview}

This chapter describes the research methodology adopted to design, implement, and validate the Behaviour-Aware Honeypot system. The methodology follows a structured, iterative engineering approach — beginning with a threat-model definition, proceeding through system design and implementation, and culminating in a multi-phase testing campaign covering unit-level component verification, scenario-based integration testing, adversarial red-team evaluation, and performance benchmarking.

The overarching research question driving the work is: \textit{can a honeypot system leverage real-time machine-learning classification and large-language-model-generated deception responses to meaningfully increase attacker dwell time, improve threat attribution accuracy, and provide actionable forensic intelligence — without incurring latency penalties that would reveal the system's artificial nature?}

\section{Research Design and Approach}

The project follows a \textbf{design science research} paradigm, in which the artefact (the honeypot system) and its evaluation are developed in tandem. Each design decision is followed by an empirical test to confirm the desired behaviour, and failing tests drive targeted refinements before the next iteration begins.

Three guiding principles governed all design choices:

\begin{enumerate}
    \item \textbf{Non-disclosure:} No response, timing artefact, or system-level fingerprint should reveal to a remote attacker that the environment is simulated.
    \item \textbf{Graceful degradation:} Every subsystem must fall back cleanly if an upstream dependency (LLM, ML model, Docker daemon) is unavailable, so the honeypot always produces a plausible response.
    \item \textbf{Non-blocking I/O:} Classification and logging must never add measurable latency to the interactive shell session, as real-world servers respond to commands in under 100\,ms.
\end{enumerate}

\section{System Architecture Overview}

The system comprises five co-operating services, each running as an isolated Docker container connected over a dedicated bridge network:

\begin{enumerate}
    \item \textbf{Core Service} — a Flask REST API that acts as the central intelligence hub. Every event generated by the trap services is forwarded here for enrichment via the ML pipeline before being returned to the caller as a structured decision payload. The core also hosts the Orchestrator, which manages the lifecycle of child containers, and the Control API, which exposes dashboard-level toggle endpoints.

    \item \textbf{SSH Honeypot} — a full Paramiko-based \ac{SSH} server bound on port 2222. It presents a realistic message-of-the-day, accepts any credential, and hands the attacker a fake shell environment backed by a static in-memory filesystem. Commands marked as attack vectors are forwarded synchronously to the Core API; all other commands are handled deterministically with zero LLM involvement.

    \item \textbf{HTTP Honeypot} — a Flask application presenting itself as the fictional company \textit{NovaTech Solutions}, running on port 8080 with forged server headers. It exposes a set of deliberately vulnerable endpoints (login form, \texttt{/.env}, \texttt{/backup}, \texttt{/admin}, \texttt{/robots.txt}) containing bait credentials that form the multi-stage kill chain.

    \item \textbf{Dashboard} — a Streamlit application providing a real-time operator view: live event feed, per-IP risk scores, behaviour state machine status, kill-chain progress, and AI-mode toggle controls.

    \item \textbf{Tunnel Services} — lightweight containers that expose the HTTP trap and dashboard externally via Cloudflare TryCloudflare, and the SSH trap via a Pinggy TCP tunnel, allowing remote attackers to reach the system over the public internet.
\end{enumerate}

\subsection{Event Life Cycle}

The following describes the path of a single attacker interaction from ingress to logged decision:

\begin{enumerate}
    \item The attacker sends an \ac{SSH} command or \ac{HTTP} request to the corresponding trap container.
    \item The trap service classifies the input locally using a rule-set and dispatches a structured JSON event to the Core API.
    \item The Core API passes the payload through the four-stage ML pipeline: TF-IDF Random Forest fast-path classifier, DistilBART zero-shot deep classifier, BiLSTM sequence behaviour profiler, and TTP extractor.
    \item The Behaviour Classifier updates the per-IP state (risk score, service set, kill-chain flags) and runs the deterministic state machine transition.
    \item If the response action is \texttt{CONTAIN}, the Maneuvering Engine generates an LLM deception response via Phi-3-mini or falls back to a rich static payload.
    \item The Core API returns the enriched decision to the trap, which renders the response to the attacker.
    \item Concurrently, the Central Logger enqueues the full event for background SHA-256 hash-chained forensic logging.
\end{enumerate}

\section{Machine Learning Pipeline}

\subsection{Fast-Path Classifier (TF-IDF + Random Forest)}

The primary classifier is a scikit-learn pipeline that vectorises each payload with a term frequency–inverse document frequency (TF-IDF) transformer and classifies the resulting feature vector using a Random Forest model trained on a labelled corpus of attack payloads. The model is loaded from disk at startup and its probability output returns a confidence distribution over seven canonical labels:

\begin{itemize}
    \item SQL Injection
    \item Command Injection
    \item Path Traversal
    \item Brute Force
    \item Credential Access
    \item Reconnaissance
    \item Benign
\end{itemize}

The fast path executes synchronously and returns in under 2\,ms, ensuring no perceptible delay is added to the attacker's interactive session.

\subsection{Deep Zero-Shot Classifier (DistilBART MNLI)}

For payloads flagged as potentially malicious by the fast-path classifier, a secondary asynchronous analysis is submitted to a thread pool. The \texttt{valhalla/distilbart-mnli-12-1} model performs zero-shot classification against the same seven-label set. Results are stored in an LRU cache and returned on the next request for the same payload, ensuring the attacker never waits for transformer inference.

A pre-processing step decodes common obfuscation layers before analysis, including hex escape sequences, Base64 strings, \ac{SQL} comment tokens, and shell token-distancing characters. The decoder iterates up to five times to handle multi-layer obfuscation.

\subsection{Behaviour Profiler (BiLSTM)}

The Behavioural BiLSTM is a two-layer, bidirectional LSTM with a hidden dimension of 128 units per direction. Its input is a matrix of DistilRoBERTa sentence embeddings of the attacker's five most recent commands. The output layer maps the hidden state to one of three attacker archetypes:

\begin{itemize}
    \item \textbf{SCRIPT\_BOT} — automated tooling with repetitive, pattern-driven commands.
    \item \textbf{PERSISTENT\_ATTACKER} — methodical manual exploration across multiple services.
    \item \textbf{APT} — low-and-slow, multi-phase adversary exhibiting kill-chain behaviour.
\end{itemize}

This classification runs exclusively in a background thread per IP. The Behaviour Classifier reads the last cached result, so the live shell response is never delayed by sequence embedding computation. When the model has not yet been trained, the system falls back gracefully with the label \texttt{UNKNOWN (Needs Training)}.

\subsection{TTP Extractor (MITRE ATT\&CK Mapping)}

The TTP Extractor scans the full command history of each session against a 12-rule MITRE ATT\&CK dictionary covering the most common honeypot-relevant techniques. The table below summarises the techniques mapped by the extractor.

\begin{table}[h!]
    \centering
    \caption{MITRE ATT\&CK Techniques Mapped by the TTP Extractor}
    \renewcommand{\arraystretch}{1.3}
    \begin{tabular}{|l|l|l|l|}
        \hline
        \textbf{ID} & \textbf{Technique} & \textbf{Tactic} & \textbf{Example Trigger} \\
        \hline
        T1110 & Brute Force & Credential Access & \texttt{hydra}, \texttt{ssh login} \\
        T1059 & Command Interpreter & Execution & \texttt{bash}, \texttt{/bin/sh}, \texttt{eval()} \\
        T1190 & Exploit Public-Facing App & Initial Access & \texttt{union select}, \texttt{' OR '1'='1} \\
        T1083 & File/Directory Discovery & Discovery & \texttt{ls -la}, \texttt{find /} \\
        T1082 & System Info Discovery & Discovery & \texttt{uname}, \texttt{whoami}, \texttt{id} \\
        T1005 & Data from Local System & Collection & \texttt{cat /etc/shadow}, \texttt{.env} \\
        T1048 & Exfiltration Alt. Protocol & Exfiltration & \texttt{curl}, \texttt{wget}, \texttt{nc} \\
        T1053 & Scheduled Task/Job & Persistence & \texttt{crontab}, \texttt{systemctl} \\
        T1136 & Create Account & Persistence & \texttt{useradd}, \texttt{adduser} \\
        T1055 & Process Injection & Defense Evasion & \texttt{LD\_PRELOAD}, \texttt{ptrace} \\
        T1014 & Rootkit & Defense Evasion & \texttt{chkrootkit}, \texttt{rkhunter} \\
        T1071 & C2 Application Protocol & Command \& Control & \texttt{http://}, \texttt{beacon} \\
        \hline
    \end{tabular}
\end{table}

The extractor assigns a threat level (BENIGN / LOW / MEDIUM / HIGH / CRITICAL) based on the number of unique techniques and the presence of high-severity tactics such as Persistence and Command-and-Control.

\subsection{Self-Healing ML Feedback Loop}

When the state machine transitions an IP to \texttt{MALICIOUS}, \texttt{CONFIRMED\_ATTACK}, or \texttt{KILL\_CHAIN\_CONFIRMED}, the classifier appends the full command history to a feedback dataset and kicks off a background retraining run in a daemon thread. This allows the Random Forest model to continuously incorporate novel attack patterns observed in the wild, without operator intervention.

\section{Deception Engine Design}

\subsection{Tiered AI Modes}

The system supports three operator-selectable AI modes, toggled live via the dashboard without restarting any container:

\begin{itemize}
    \item \textbf{Tier 1 — Eco:} The LLM and BiLSTM are bypassed entirely. The fast-path rule engine and static bait payloads handle all responses. Suitable for low-resource environments or initial deployment.
    \item \textbf{Tier 2 — Standard:} The BiLSTM profiler runs asynchronously and the fast-path Random Forest classifier is active. Static bait payloads are used for terminal output. Suitable for production with limited resources.
    \item \textbf{Tier 3 — Advanced:} Full pipeline active. Attack-class commands trigger real-time Phi-3-mini generation via the Ollama local inference server. Suitable for high-fidelity deception and extended attacker engagement.
\end{itemize}

\subsection{LLM Response Generation and Sanitisation}

In Tier 3 mode, the Maneuvering Engine builds a per-call system prompt from the attacker's current session state (hostname, username, current working directory, files visible in that directory), then submits the attacker's command to Phi-3-mini via the Ollama \ac{HTTP} API with a 3-second deadline.

The raw LLM output passes through a sanitisation pipeline before delivery:

\begin{enumerate}
    \item \textbf{Markdown fence stripping} — responses accidentally wrapped in code fences are unwrapped.
    \item \textbf{Hostname correction} — known-bad hostnames are replaced with the session's seeded hostname.
    \item \textbf{Corruption signal rejection} — phrases such as ``this is a honeypot'' or ``I am an AI'' cause the entire response to be discarded and replaced with a static fallback.
    \item \textbf{Leaked prompt stripping} — a trailing shell prompt hallucinated by the LLM is removed.
    \item \textbf{Sensitive-content gating} — responses containing credential markers for commands that did not explicitly request file reads are rejected, preventing unsolicited secret-leaking.
\end{enumerate}

Responses are cached in an LRU cache keyed on the command, current working directory, hostname, username, and available files, so repeat commands from the same session produce identical output — a critical consistency requirement.

\subsection{Adaptive Response Timing}

The engine introduces a small, attacker-class-dependent artificial delay before each response to mimic realistic server processing time without being detectable.

\begin{table}[h!]
    \centering
    \caption{Adaptive Delay by Attacker Class}
    \renewcommand{\arraystretch}{1.3}
    \begin{tabular}{|l|l|}
        \hline
        \textbf{Attacker Class} & \textbf{Injected Delay} \\
        \hline
        SCRIPT\_BOT & 400\,ms \\
        PERSISTENT\_ATTACKER & 150\,ms \\
        APT & 50\,ms \\
        \hline
    \end{tabular}
\end{table}

\subsection{Deterministic Command Boundary}

To prevent LLM involvement from introducing inconsistency for simple shell interactions, a strict deterministic command boundary is enforced. Navigation commands (\texttt{cd}, \texttt{pwd}, \texttt{ls}), identity commands (\texttt{whoami}, \texttt{id}, \texttt{hostname}), and all file-reading commands (\texttt{cat}) are always handled by the local fake shell with zero LLM calls. The canonical file-contents dictionary is the single source of truth; the LLM may never contradict it.

\section{Kill Chain Design}

A deliberate multi-stage kill chain is seeded across both services to attract and track sophisticated attackers. The chain proceeds as follows:

\begin{enumerate}
    \item The attacker visits the \ac{HTTP} trap and reads \texttt{/robots.txt}, which discloses sensitive paths (\texttt{/backup}, \texttt{/.env}, \texttt{/admin}).
    \item The attacker retrieves \texttt{/.env}, which contains a working \ac{SSH} password for a privileged user on the trap server.
    \item The attacker connects to the \ac{SSH} trap on port 2222 with those credentials.
    \item The \ac{SSH} service recognises the kill-chain credential pair and fires the \texttt{SSH\_KILL\_CHAIN\_LOGIN} event.
    \item The Behaviour Classifier immediately transitions the IP to \texttt{KILL\_CHAIN\_CONFIRMED}, raising the risk score and marking the kill chain complete.
    \item Optionally, the Orchestrator spins up a dedicated, unprivileged Alpine container for physical sandboxing of the attacker IP, providing true isolation beyond LLM spoofing.
\end{enumerate}

\subsection{HTTP Trap Endpoints}

The table below summarises the \ac{HTTP} trap's deception endpoints and their associated MITRE technique mappings.

\begin{table}[h!]
    \centering
    \caption{HTTP Honeypot Endpoints and MITRE ATT\&CK Mapping}
    \renewcommand{\arraystretch}{1.3}
    \begin{tabular}{|l|l|l|}
        \hline
        \textbf{Endpoint} & \textbf{Deception Payload} & \textbf{MITRE ID} \\
        \hline
        \texttt{/robots.txt} & Discloses \texttt{/admin}, \texttt{/backup}, \texttt{/.env} & T1595.002 \\
        \texttt{/.env} & SSH credentials, AWS keys, DB password & T1552.001 \\
        \texttt{/backup/server\_info.txt} & SSH host, port, and user details & T1552.001 \\
        \texttt{/backup/db\_backup\_nov.sql} & Fake user table with bcrypt hashes & T1552.001 \\
        \texttt{/login} & SQLi/XSS detection and fake SQL error & T1190, T1059.007 \\
        \texttt{/search} & SQLi/XSS/Path Traversal detection & T1083, T1190 \\
        \texttt{/download} & Path traversal detection & T1083 \\
        \texttt{/admin} & 403 with event logged & T1595.002 \\
        \texttt{/api/users} & 401 with DB schema leak in debug field & T1595.002 \\
        \hline
    \end{tabular}
\end{table}

\section{Forensic Logging Design}

The Central Logger implements a non-blocking, hash-chained forensic event log. The public \texttt{log(event)} method appends an event to an unbounded queue and sets a threading event; a single daemon worker thread drains the queue sequentially. This design prevents any Flask request handler thread from ever blocking on disk I/O.

Each log entry is enriched with a SHA-256 hash of its own content and the hash of the previous entry, forming a blockchain-style chain. The genesis block uses a constant string as the initial previous hash. This chain construction ensures that post-hoc tampering with any historical event is detectable by chain validation — a key forensic integrity requirement.

\section{Testing Methodology}

Testing was structured across five distinct phases, arranged in increasing complexity and realism:

\begin{enumerate}
    \item \textbf{Phase 0} — Environment and health verification.
    \item \textbf{Phase 1} — Control plane and toggle validation.
    \item \textbf{Phase 2} — Normal-activity baseline testing.
    \item \textbf{Phase 3} — Output consistency testing.
    \item \textbf{Phase 4} — Adversarial attack simulation.
    \item \textbf{Phase 5} — Forensic log integrity check.
\end{enumerate}

\subsection{Phase 0 — Environment and Health Verification}

Before any functional testing, an automated validation script confirms that all critical infrastructure is healthy by verifying the following:

\begin{itemize}
    \item The Core API status endpoint returns HTTP 200.
    \item The \ac{HTTP} and \ac{SSH} honeypot status endpoints are reachable.
    \item The tunnel status endpoint returns active Cloudflare and Pinggy URLs.
    \item The Ollama LLM server is reachable and the Phi-3-mini model is listed in the tag manifest.
\end{itemize}

\subsection{Phase 1 — Control Plane and Toggle Validation}

The system must support live reconfiguration without container restarts. This phase exercises the \ac{HTTP} and \ac{SSH} stop/start cycles and AI mode switching across all three tiers. Each toggle is followed by a short sleep to allow the Docker daemon to reflect the new state, then a status check to confirm the expected running or stopped condition.

\subsection{Phase 2 — Normal Activity Baseline}

A set of representative legitimate shell commands is executed via a Paramiko connection to verify that the fake shell handles all deterministic command classes correctly and within an acceptable latency budget. The commands tested include \texttt{whoami}, \texttt{pwd}, \texttt{ls -la}, \texttt{cd /var/log}, and \texttt{cat /etc/os-release}. The expected output for each command is validated against the canonical fake filesystem and file-contents data structures. Latency for all deterministic commands must not exceed 100\,ms.

\subsection{Phase 3 — Consistency Testing}

A core non-disclosure requirement is output idempotency: repeated identical commands from the same session must always return byte-for-byte identical output. This phase runs \texttt{ls} and \texttt{cat /etc/passwd} three times each on the same shell session and asserts all three outputs are equal.

Session-to-session consistency of random metadata (file sizes, timestamps, system load) is guaranteed by seeding the per-session random number generator with a hash of the client IP, so the same attacker always sees the same filesystem metadata across reconnections.

\subsection{Phase 4 — Adversarial Attack Simulation}

\subsubsection{SSH Attack Vectors}

The following \ac{SSH}-layer payloads are tested to verify correct detection, classification, and deception routing:

\begin{table}[h!]
    \centering
    \caption{SSH Attack Vectors Tested}
    \renewcommand{\arraystretch}{1.3}
    \begin{tabular}{|l|l|l|}
        \hline
        \textbf{Payload} & \textbf{Expected Classification} & \textbf{Expected Action} \\
        \hline
        \texttt{; ls} & Command Injection (ATTACK) & LLM or static bait \\
        \texttt{\&\& whoami} & Command Injection (ATTACK) & LLM or static bait \\
        \texttt{bash -i >\& /dev/tcp/...} & Command Injection (ATTACK) & LLM or static bait \\
        \texttt{echo d2hvYW1p | base64 -d} & Command Injection (ATTACK) & LLM or static bait \\
        \texttt{cat /etc/passwd} & CONTENT (deterministic) & Fake file contents \\
        \texttt{find / -perm -4000} & CONTENT (deterministic) & Bait SUID list \\
        \texttt{sudo -l} & Privilege escalation event & Fake password prompt \\
        \texttt{su root} & SYSTEM (deterministic) & Fake prompt \\
        \texttt{cat ../../../etc/shadow} & CONTENT (deterministic) & Permission denied \\
        \hline
    \end{tabular}
\end{table}

\subsubsection{HTTP Attack Vectors}

The following \ac{HTTP} payloads are posted to the \texttt{/login} endpoint to verify pattern detection:

\begin{table}[h!]
    \centering
    \caption{HTTP Attack Vectors Tested}
    \renewcommand{\arraystretch}{1.3}
    \begin{tabular}{|l|l|l|}
        \hline
        \textbf{Payload} & \textbf{Expected Detection} & \textbf{Expected Response} \\
        \hline
        \texttt{' OR 1=1 --} & SQL Injection & Fake SQL error message \\
        \texttt{' UNION SELECT NULL,NULL--} & SQL Injection & Fake SQL error message \\
        \texttt{../../../../etc/passwd} & Path Traversal & 400/403 \\
        \texttt{; whoami} & Command Injection & HTTP 200 (event logged) \\
        \texttt{| id} & Command Injection & HTTP 200 (event logged) \\
        Oversized payload (500+ chars) & Fuzzing & HTTP 200 (event logged) \\
        \hline
    \end{tabular}
\end{table}

\subsubsection{Obfuscated Payload Handling}

A specific sub-test verifies the pre-processing pipeline's ability to decode obfuscated payloads before classification:

\begin{itemize}
    \item \textbf{Hex-encoded commands:} A hex-encoded payload decodes to \texttt{rm -rf} prior to classification.
    \item \textbf{Base64-encoded commands:} A Base64 payload is decoded so the underlying \texttt{whoami} intent is correctly labelled Reconnaissance.
    \item \textbf{Shell token distancing:} Shell metacharacter tokens are stripped before analysis to reveal the true command.
\end{itemize}

\subsubsection{Kill Chain End-to-End Test}

A dedicated integration test traverses the full kill chain in sequence:

\begin{enumerate}
    \item \texttt{GET /robots.txt} — verifies disclosure of sensitive paths and that the corresponding event is logged.
    \item \texttt{GET /.env} — verifies \ac{SSH} credential exposure and that the kill-chain accessed flag is set.
    \item \ac{SSH} connect with the kill-chain credentials — verifies that the kill-chain login event is fired.
    \item Verify \texttt{behaviour\_state = KILL\_CHAIN\_CONFIRMED} and \texttt{kill\_chain.complete = true} in the Core API response.
    \item Verify that the risk score exceeds the expected threshold.
\end{enumerate}

\subsection{Phase 5 — Log Integrity Verification}

After all attack simulations, the forensic event log is opened and its hash chain is validated:

\begin{enumerate}
    \item The first entry must have a previous hash equal to the genesis block constant.
    \item For every subsequent entry, the SHA-256 hash of the previous entry is recomputed and confirmed to match the current entry's previous-hash field.
    \item The total event count is confirmed to match the number of test interactions performed.
\end{enumerate}

\section{Scenario-Based Attacker Profile Testing}

In addition to the structured validation phases, four distinct attacker profiles were simulated end-to-end to evaluate the behaviour state machine and risk-scoring logic under realistic event sequences. Each profile was assigned a unique source IP and injected events directly via the Core API, providing deterministic and reproducible execution.

\subsection{Profile 1 — Dumb Bot}

This profile simulates an automated scanner with no intelligence, issuing GET requests to common paths followed by credential brute forcing. The results are summarised in the table below.

\begin{table}[h!]
    \centering
    \caption{Dumb Bot Scenario Results}
    \renewcommand{\arraystretch}{1.3}
    \begin{tabular}{|c|l|l|l|l|l|}
        \hline
        \textbf{Step} & \textbf{Action} & \textbf{Attack Type} & \textbf{State} & \textbf{Risk} & \textbf{MITRE} \\
        \hline
        1 & \texttt{/robots.txt} & BENIGN & NEW & 0.0 & T1071 \\
        2 & \texttt{/wp-login.php} & BENIGN & PROBING & 0.0 & T1071, T1110 \\
        3 & \texttt{/.env} & CRED\_ACCESS & SUSPICIOUS & 3.5 & T1071, T1110, T1005 \\
        4 & \texttt{admin:admin} & BRUTE\_FORCE & MALICIOUS & 10.5 & T1071, T1110, T1005 \\
        5 & \texttt{root:root} & BRUTE\_FORCE & MALICIOUS & 19.0 & T1071, T1110, T1005 \\
        \hline
    \end{tabular}
\end{table}

The system correctly escalated from NEW through PROBING and SUSPICIOUS to MALICIOUS, mapping three distinct MITRE techniques across the session and triggering CONTAIN actions from Step 3 onward.

\subsection{Profile 2 — Script Kiddie}

This profile simulates a moderately skilled attacker using publicly available tools: an \ac{SQL} injection payload, a path traversal attempt, a command injection probe, a brute-force login, and a malware download attempt.

The system correctly identified SQL Injection immediately on the first event and escalated to MALICIOUS by Step 2. By Step 5, six distinct MITRE techniques had been mapped (T1190, T1110, T1083, T1059, T1048, T1071) and the risk score had reached 30.0.

\subsection{Profile 3 — Persistent Attacker (Kill Chain)}

This profile is the primary kill-chain scenario, designed to trigger the full credential-to-shell path:

\begin{itemize}
    \item \textbf{Steps 1–2:} Benign reconnaissance. State remains NEW then PROBING.
    \item \textbf{Step 3:} Backup file access event. Risk rises to 3.0; AI classification updates to PERSISTENT\_ATTACKER.
    \item \textbf{Step 4:} Kill-chain \ac{SSH} login with bait credentials. State immediately transitions to \texttt{KILL\_CHAIN\_CONFIRMED}; risk jumps to 18.5.
    \item \textbf{Steps 5–6:} Post-compromise enumeration accumulates further MITRE techniques; risk reaches 34.0.
\end{itemize}

\subsection{Profile 4 — APT Stealth}

This profile simulates a low-and-slow adversary that avoids triggering simple threshold-based defences. Initial \ac{HTTP} access is entirely benign. A single \ac{SSH} probe moves the state to SUSPICIOUS. Subsequent commands escalate through MALICIOUS to CONFIRMED\_ATTACK.

The critical observation is that a log-wiping command (\texttt{rm -rf /var/log/auth.log}) was correctly detected as Command Injection (T1070 — Indicator Removal on Host) even though the preceding privilege-check command was classified BENIGN, demonstrating that the risk accumulator correctly combines historical context with per-event classification.

\section{Red Team Testing}

An independent red-team test was conducted against the live system. Unlike the internal scenario scripts which inject events directly into the Core API, the red-team test operates over a real Paramiko \ac{SSH} connection and real \ac{HTTP} requests — exposing the system to the same code paths a genuine attacker would traverse.

The test covered nine scenario groups across twenty individual test cases:

\begin{enumerate}
    \item \textbf{Basic Shell Behaviour} — \texttt{pwd}, \texttt{ls}, valid and invalid \texttt{cd}, and \texttt{touch}. Verified correct navigation, error messages, and directory state.
    \item \textbf{File Access} — \texttt{cat} of existing and non-existing files. Verified correct fake-file delivery and proper error messaging for missing files.
    \item \textbf{Output Consistency} — Three consecutive \texttt{ls} calls and two consecutive \texttt{cat} calls on the same session. Verified identical byte-level output across all repetitions.
    \item \textbf{Invalid Commands} — Arbitrary strings. Verified that a command-not-found error is always returned with no LLM hallucination.
    \item \textbf{Enumeration} — \texttt{whoami} and \texttt{id}. Verified that the UID derived from the username hash was consistent between both commands.
    \item \textbf{Attack Simulation (\ac{SSH})} — Shell metacharacter injection, a TCP reverse-shell payload, and a Base64-encoded payload. Each triggered ATTACK classification and the CONTAIN response path.
    \item \textbf{\ac{HTTP} Attacks} — \ac{SQL} injection, path traversal, and command injection via \texttt{POST /login}. All detected and logged correctly.
    \item \textbf{Context Consistency} — Navigation to a subdirectory followed by \texttt{pwd} and \texttt{ls}. Verified that working directory state was maintained across commands within a session.
    \item \textbf{Performance} — Back-to-back \texttt{ls} and \texttt{pwd}. Verified sub-100\,ms response times for all deterministic commands.
\end{enumerate}

\section{Performance Requirements and Results}

Three latency targets were defined prior to implementation, derived from the non-disclosure requirement that the system must not be distinguishable from a legitimate server by response-time analysis.

\begin{table}[h!]
    \centering
    \caption{Performance Targets and Measurements}
    \renewcommand{\arraystretch}{1.3}
    \begin{tabular}{|l|l|l|l|}
        \hline
        \textbf{Category} & \textbf{Target} & \textbf{Measured (p95)} & \textbf{Result} \\
        \hline
        Deterministic shell commands & $<100$\,ms & $\approx 15$\,ms & PASS \\
        ML classification (fast-path) & $<2$\,ms & $<2$\,ms & PASS \\
        LLM response (Tier 3, Phi-3-mini) & $<3$\,s & $\approx 1.2$–$2.8$\,s & PASS \\
        Event logging (async, non-blocking) & 0\,ms added to handler & 0\,ms & PASS \\
        DistilBART deep analysis & Non-blocking (async) & $\approx 800$\,ms (background) & PASS \\
        \hline
    \end{tabular}
\end{table}

All deterministic shell commands returned within 15\,ms on average, well below the 100\,ms threshold. LLM-generated responses carried an additional 1.2–2.8 seconds of latency, which was judged acceptable because attack tools do not enforce sub-second response deadlines, and the adaptive delay heuristics contextualise this as normal server processing time.

\section{Limitations and Mitigations}

\paragraph{LLM Identity Drift}
The Phi-3-mini model occasionally generates responses that reference incorrect hostnames or identity strings inconsistent with the session context. This is mitigated by the multi-layer sanitisation pipeline described in Section~\ref{sec:llm_sanitation}. Responses that cannot be corrected are silently discarded and a static fallback payload is returned instead, ensuring the honeypot never returns an empty response.

\paragraph{BiLSTM Cold Start}
The BiLSTM model returns \texttt{UNKNOWN (Needs Training)} until a sufficient corpus of labelled sessions has been accumulated and the model retrained. During this period, the Behaviour Classifier falls back to the rule-based state machine for attacker-class routing, which remains effective for all core response-selection logic.

\paragraph{Tunnel Reliability}
External accessibility depends on third-party tunnel services (Cloudflare TryCloudflare, Pinggy). Both provide fallback URL patterns; the Orchestrator also exposes \texttt{localhost:2222} as a permanent \ac{SSH} fallback, ensuring local testing and demonstrations remain unaffected by tunnel outages.

\paragraph{Single-Host Deployment}
The current implementation runs all containers on a single Docker host. In a production setting, individual services could be distributed across separate hosts for higher throughput and to prevent a compromised container from affecting the Core API. The Docker socket mounting required by the Orchestrator represents an acknowledged privilege escalation risk on the host, which should be addressed via rootless Docker in future iterations.

\section{Summary}

This chapter presented the complete methodology and testing framework for the Behaviour-Aware Honeypot system. The design science approach guided iterative development, with each subsystem — the four-stage ML pipeline, the tiered deception engine, the kill-chain infrastructure, and the forensic logger — validated through increasingly realistic test scenarios.

The five-phase test plan confirmed that all core requirements were met: low-latency deterministic responses, correct attack classification and MITRE mapping, stateful behaviour escalation, kill-chain detection, idempotent output for identical commands, and tamper-evident hash-chained logging. Scenario testing across four attacker profiles demonstrated accurate progression through the state machine from NEW to KILL\_CHAIN\_CONFIRMED, with risk scores and MITRE attribution accumulating consistently with observed behaviour.
